\section{Reconstructing a reflection from two chamber orbits}
\label{sec:chambers}
A type-preserving map of chambers preserves an $i$-adjacency when it preserves the operation of replacing the type-$i$ object in a chamber. In a thin geometry this operation has a unique value different from the original chamber.

\begin{lemma}[Propagation and reconstruction]\label{lem:propagation}
In a residually connected incidence geometry, the chambers containing any fixed flag are connected by adjacency steps that retain that flag. If the geometry is thin, a type-preserving automorphism fixing a chamber is the identity. Moreover, every permutation of chambers preserving every adjacency color comes from a unique type-preserving automorphism of the geometry.
\end{lemma}
\begin{proof}
We first prove chamber connectedness by induction on the rank of a residue. A residue of rank zero has one chamber. In rank one any two chambers are adjacent. At higher rank its incidence graph is connected. Join an object in the first chamber to an object in the second by an incidence path. Consecutive objects lie together in a chamber, since every flag extends. Successive chambers chosen in this way share an object; the induction hypothesis in that object's residue joins them while retaining it. Concatenating these paths joins the original chambers. Applying the same argument inside the residue of a fixed flag proves the first assertion.

In a thin geometry a chamber-fixing automorphism fixes the unique adjacent chamber of each color. Induction along a chamber path fixes every chamber and hence every object. For reconstruction, the chambers containing a type-$i$ object form exactly one connected component after all $i$-edges are removed. One inclusion follows because the remaining moves preserve that object; the opposite inclusion is the first assertion. A color-preserving chamber permutation therefore permutes these components, inducing a bijection on the objects of each type. Two objects are incident exactly when they occur in a common chamber, again by flag extension. Thus the induced map preserves incidence in both directions. Its chamber action is the prescribed one, so it is unique.
\end{proof}

\begin{theorem}[Two regular orbits]\label{thm:reflection}
Let $G=\langle\alpha_i:i\in I\rangle$, where $0\in I$ and $\alpha_0=1$. Suppose a thin, residually connected geometry has chambers
\[
             (g,+),\ (g,-)\qquad(g\in G)
\]
with $G$ acting by left multiplication in each copy, and with its $i$-adjacencies given by
\begin{equation}\label{eq:edge}
                   (g,+)\ \sim_i\ (g\alpha_i^{-1},-).
\end{equation}
Then its full type-preserving automorphism group is $G$ unless an automorphism $\theta\in\Aut(G)$ satisfies
\begin{equation}\label{eq:inverse}
                       \theta(\alpha_i)=\alpha_i^{-1}
                       \quad(i\in I).
\end{equation}
If such $\theta$ exists, the full group is $G\rtimes\langle F\rangle$, where $F^2=1$, $FgF^{-1}=\theta(g)$, and
\[
 F(g,+)=(\theta(g),-),\qquad F(g,-)=(\theta(g),+).
\]
In the first case the geometry is chiral; in the second it is regular.
\end{theorem}
\begin{proof}
An automorphism satisfying \eqref{eq:inverse} squares to the identity on every generator, and hence on $G$. The displayed map $F$ is therefore an involution. An edge in \eqref{eq:edge} maps to
\[
        (\theta(g),-)\ \sim_i\ (\theta(g)\alpha_i,+),
\]
which is an edge by the same formula. Lemma~\ref{lem:propagation} extends $F$ to an automorphism of the geometry. Direct evaluation on both chamber copies gives $FgF^{-1}=\theta(g)$.

Conversely, let $A$ be the full type-preserving automorphism group. It acts freely on chambers by Lemma~\ref{lem:propagation}. Evaluation at $(1,+)$ identifies its elements whose image lies in the positive copy exactly with $G$. If $A$ also reaches the negative copy, compose one such element with a member of $G$ to obtain $F$ taking $(1,+)$ to $(1,-)$. This normalized element is unique. Since $0$-adjacency exchanges these two chambers, $F$ also takes $(1,-)$ to $(1,+)$; freeness gives $F^2=1$. There are exactly two cosets of $G$ in $A$, so $G$ is normal.

Write $\theta(g)=FgF^{-1}$. The chamber $(\alpha_i,+)$ is obtained from $(1,+)$ by a $0$-step and then an $i$-step. Applying $F$ gives the $i$-neighbor of $(1,+)$, namely $(\alpha_i^{-1},-)$. On the other hand $F(\alpha_i,+)=\theta(\alpha_i)(1,-)$. Freeness on the negative copy yields \eqref{eq:inverse}. Thus orbit fusion is equivalent to the specified inversion, and without it the full group has exactly the two alternating chamber orbits. No finite cardinality assumption entered the argument.
\end{proof}

This is an explicit version of the marked-inversion criterion in Fernandes--Leemans--Weiss, \emph{Highly symmetric hypertopes}, Theorem 8.2. The formula for $F$ discharges the extension from a group automorphism to an incidence automorphism. It does not identify an arbitrary change of matrix coordinates with a geometric reflection.
